% =====================================================================
%  deck_kernels.tex  --  meeting deck, 2026-09-04 (beamer, 16:9)
%
%  Framing: reading week. The supervisor's two pointers, the three
%  papers in related_works/, and where they would fit our model.
%  NO measured results appear here (those are RESERVE material in
%  docs/meeting_qa_prep.md, deployed only if Ray chooses in the room).
%
%  Style: house rules from the July lab talk. No em-dashes, no
%  semicolons, no asterisks on slides. One term per concept. Symbols
%  defined with role and units. Only show what is used.
%
%  Compile:  cd docs/slides && tectonic deck_kernels.tex
% =====================================================================
\documentclass[aspectratio=169,11pt]{beamer}
\usetheme{default}
\usefonttheme{professionalfonts}
\setbeamertemplate{navigation symbols}{}
\setbeamertemplate{itemize items}{\textcolor{okblue}{\scriptsize$\blacksquare$}}
\setbeamertemplate{footline}{%
  \hfill\scriptsize\textcolor{muted}{\insertframenumber\,/\,\inserttotalframenumber}\hspace{2mm}\vspace{2mm}}

\usepackage{amsmath,amssymb,bm}
\usepackage{booktabs}
\usepackage{graphicx}
\usepackage{tikz}

\definecolor{okblue}{HTML}{0072B2}
\definecolor{okverm}{HTML}{D55E00}
\definecolor{okgreen}{HTML}{009E73}
\definecolor{okamber}{HTML}{E69F00}
\definecolor{muted}{HTML}{6B7280}
\definecolor{ink}{HTML}{1A1A1A}
\setbeamercolor{frametitle}{fg=ink}
\setbeamercolor{normal text}{fg=ink}
\setbeamerfont{frametitle}{series=\bfseries,size=\large}

\title{\textbf{Reading week: spectral kernels for pitch curves}}
\subtitle{Three papers, and where they fit the score graph model}
\author{Raynaldi Lalang}
\institute{Kyoto University}
\date{September 4, 2026}

\begin{document}

\begin{frame}
  \titlepage
\end{frame}

% ---------------------------------------------------------------------
\begin{frame}{Why these papers}
  \textbf{Your two pointers last time:}
  \begin{itemize}
    \item periodic kernel Gaussian processes
    \item generalized spectral mixtures
  \end{itemize}
  \vspace{2mm}
  \textbf{The question they answer.} Our model summarizes a note's pitch
  curve with one parametric shape,
  \[
    \mathrm{cents}(t) \;=\; c \;+\;
    \mathbf{1}[t \ge \delta]\,\gamma \sin\!\bigl(2\pi f (t-\delta)\bigr)
  \]
  \begin{itemize}
    \item $c$ is the pitch centre in cents, $\gamma$ the vibrato extent
          in cents, $f$ the vibrato rate in Hz, $\delta$ the onset delay
          in seconds
    \item every parameter is held constant across the note
    \item the papers ask instead: what if the curve itself gets a
          Gaussian process prior, and the kernel carries the structure
  \end{itemize}
\end{frame}

% ---------------------------------------------------------------------
\begin{frame}{Paper 1. Wilson and Adams 2013, the spectral mixture kernel}
  \textbf{Idea.} Do not pick a kernel. Model its spectrum as a Gaussian
  mixture, the kernel follows in closed form,
  \[
    k(\tau) \;=\; \sum_{q=1}^{Q} w_q\,
      \exp\!\bigl(-2\pi^2 \tau^2 v_q\bigr)\,
      \cos\!\bigl(2\pi \tau \mu_q\bigr)
  \]
  \begin{itemize}
    \item $\mu_q$ is a frequency in Hz, $v_q$ its spread, so
          $1/\sqrt{v_q}$ is how long the oscillation stays coherent,
          $w_q$ its weight
    \item enough components approximate any stationary kernel
    \item the marginal likelihood prunes components that the data does
          not need
    \item their demos extrapolate CO$_2$ and airline data where the
          squared exponential kernel only interpolates
  \end{itemize}
  \vspace{1mm}
  \textbf{Fit to us.} A note's vibrato is one component at $\mu \approx f$
  with finite coherence. Drift is a second component near zero frequency.
  Vibrato plus drift is a two component spectral mixture prior on the
  cents curve.
\end{frame}

% ---------------------------------------------------------------------
\begin{frame}{Paper 2. Remes, Heinonen, Kaski 2017, the generalized spectral mixture}
  \textbf{Idea.} Promote the mixture constants to functions of time,
  each with its own Gaussian process prior,
  \[
    \log w_i(t), \quad \log \ell_i(t), \quad \operatorname{logit} \mu_i(t)
    \;\sim\; \mathcal{GP}
  \]
  \begin{itemize}
    \item frequency, weight and coherence length may all drift smoothly
    \item the stationary spectral mixture is the constant function
          special case
    \item their first experiment recovers an oscillation whose frequency
          changes over time
  \end{itemize}
  \vspace{1mm}
  \textbf{Fit to us.} Vibrato rate that wanders and extent that grows
  within one note, as smooth functions. Our sine fit is the degenerate
  limit, one component with constant functions and a hard onset gate.
  \vspace{1mm}

  \textbf{Honest cost.} Hyperparameters become latent functions, so
  inference is a heavier optimization. On a short note a free frequency
  function can imitate a drift, which is the confound our hard
  identifiability rules currently manage.
\end{frame}

% ---------------------------------------------------------------------
\begin{frame}{Paper 3. Alvarado and Stowell 2016, Gaussian processes for music audio}
  \textbf{Idea.} A recording is a sum of per note processes gated by
  smooth on and off windows. Each note's kernel is an exponentiated
  cosine,
  \[
    k_{\mathrm{EC}}(\tau) \;=\; \sigma^2 \exp\!\bigl(z \cos(\omega \tau)\bigr)
  \]
  \begin{itemize}
    \item $\omega$ is the fundamental in rad/s and $z$ the spectral
          richness, the spectrum is a harmonic comb, fundamental plus
          partials
    \item multiplying by a squared exponential term lets the envelope
          move
    \item pitch estimation slides $\omega$ and reads the marginal
          likelihood
  \end{itemize}
  \vspace{1mm}
  \textbf{Fit to us.} They must hand specify every note window. The
  score gives us onsets and durations for free. Their evidence scan over
  $\omega$ is the same inference pattern as our collapsed audio
  likelihood.
\end{frame}

% ---------------------------------------------------------------------
\begin{frame}{One story, and one deliberate difference}
  \textbf{The ladder.}
  \begin{itemize}
    \item Wilson and Adams give the right prior family for a quasi
          periodic curve
    \item Remes et al.\ extend it to rates and extents that drift, which
          is what real vibrato does
    \item Alvarado and Stowell assemble the same ideas at audio level
          with per note gating
  \end{itemize}
  \vspace{2mm}
  \textbf{The design axis to be explicit about.} Kernel side
  constructions marginalize the note's structure away. We keep named per
  note variables, $c$, $\gamma$, $f$, because the variables themselves
  are the object of study, the transcription output, and the thing the
  graph couples across notes. Both designs are valid. Ours is chosen,
  not naive.
\end{frame}

% ---------------------------------------------------------------------
\begin{frame}{Where the kernels would plug in}
  \textbf{Slot A, a curve level estimator.} Gaussian process regression
  with a spectral mixture prior on each note's cents curve, replacing
  the parametric sine fit. The output is a posterior over slowly varying
  centre, extent and rate rather than point values with delta method
  variances.
  \vspace{2mm}

  \textbf{Slot B, the audio model's curve prior.} The waveform
  likelihood needs a prior over pitch curve deviations. These kernel
  families are the principled choice, they stay Gaussian, so the closed
  form machinery carries over.
  \vspace{2mm}

  \textbf{Boundary.} Both slots concern structure within a note. The
  graph prior across notes is a separate, already validated layer and is
  untouched by this.
\end{frame}

% ---------------------------------------------------------------------
\begin{frame}{What I would like from you today}
  \begin{itemize}
    \item were these the three papers you had in mind, and are there
          others you would add
    \item which slot deserves priority, the curve level estimator or
          the audio model's curve prior
    \item my own next steps either way: consolidate fundamentals with a
          structured study plan, and simplify the thesis draft so the
          main line is readable before new modelling work
  \end{itemize}
\end{frame}

\end{document}
